\Ch{Basic Concepts}{Basic}

\begin{note}
  \p HLSL inherits a significant portion of its language semantics from C and
  C++. Some of this is a result of intentional adoption of syntax early in the
  development of the language and some a side-effect of the Clang-based
  implementation of DXC.

  \p This chapter includes a lot of definitions that are inherited from C and
  C++. Some are identical to C or C++, others are slightly different. HLSL is
  neither a subset nor a superset of C or C++, and cannot be simply described in
  terms of C or C++. This specification includes all necessary definitions for
  clarity.
\end{note}

\Sec{Preamble}{Basic.preamble}

\p An \textit{entity} is a value, object, function, enumerator, type, class
member, bit-field, template, template specialization, namespace, or pack.

\p A \textit{name} is a use of an \textit{identifier} (\ref{Expr.Primary.ID}),
\textit{operator-function-id} (\ref{Overload.operator}),
\textit{conversion-function-id} (\ref{Classes.Conversions}),
or \textit{template-id} (\ref{Template}) that denotes any entity or
\textit{label} (\ref{Stmt.Label}).

\p Every name that denotes an entity is introduced by a \textit{declaration}.
Every name that denotes a label is introduced by a \textit{labeled statement}
(\ref{Stmt.Label}).

\p A \textit{variable} is introduced by the declaration of a reference other
than a non-static data member of an object. The variable's name denotes the
reference or object.

\p Whenever a name is encountered it is necessary to determine if the name
denotes an entity that is a type or template. The process for determining if a
name refers to a type or template is called \textit{name lookup}.

\p Two names are the same name if:
\begin{itemize}
\item they are identifiers comprised of the same character sequence, or
\item they are operator-function-ids formed with the same operator, or
\item they are conversion-function-ids formed with the same type, or
\item they are template-ids that refer to the same class or function.
\end{itemize}

\begin{note}
  This section matches \gls{isoCPP} section \textbf{[basic]} except for the
  exclusion of \texttt{goto} and \textit{literal operators}.
\end{note}

\Sec{Declarations and definitions}{Basic.Decl}

\p A declaration (\ref{Decl}) may introduce one or more names into a translation
unit or redeclare names introduced by previous declarations. If a declaration
introduces names, it specifies the interpretation and attributes of these names.
A declaration may also have effects such as:
\begin{itemize}
\item verifying a static assertion (\ref{Decl}),
\item use of attributes (\ref{Decl}), and
\item controlling template instantiation (\ref{Template.Inst}).
\end{itemize}

\p A declaration is a \textit{definition} unless:
\begin{itemize}
\item it declares a function without specifying the function's body
(\ref{Decl.Function}),
\item it is a parameter declaration in a function declaration that does not
specify the function's body (\ref{Decl.Function}),
\item it is a global or namespace member declaration without the \texttt{static}
specifier,
\item it declares a static data member in a class definition,
\item it is a class name declaration,
\item it is a template parameter,
\item it is a \texttt{typedef} declaration (\ref{Decl}),
\item it is an \textit{alias-declaration} (\ref{Decl}),
\item it is a \textit{using-declaration} (\ref{Decl}),
\item it is a \textit{static\_assert-declaration} (\ref{Decl}),
\item it is an \textit{empty-declaration} (\ref{Decl}),
\item or a \textit{using-directive} (\ref{Decl}).
\end{itemize}

\begin{note}
  Global variable declarations are implicitly constant and external in HLSL.
\end{note}

\p The two examples below are adapted from \gls{isoCPP} \textbf{[basic.def]}. All
but one of the following are definitions:
\begin{HLSL}
int f(int x) { return x+1; } // defines f and x
struct S {int a;int b;};     // defines S, S::a, and S::b
struct X {                   // defines X
  int x;                     // defines non-static member x
  static int y;              // declares static data member y
};
int X::y = 1;                // defines X::y
enum { up, down };           // defines up and down
namespace N {                // defines N
int d;                       // declares N::d
static int i;                // defines N::i
}
\end{HLSL}

\p All of the following are declarations:
\begin{HLSL}
int a;                       // declares a
const int c;                 // declares c
X anX;                       // declares anX
int f(int);                  // declares f
struct S;                    // declares S
typedef int Int;             // declares Int
using N::d;                  // declares d
using Float = float;         // declares Float
cbuffer CB {                 // does not declare CB
  int z;                     // declares z
}
tbuffer TB {                 // does not declare TB
  int w;                     // declares w
}
\end{HLSL}

\Sec{One-Definition Rule}{Basic.ODR}

\p The \gls{isoCPP} \textit{One-definition rule} is adopted as defined in
\gls{isoCPP} \textbf{[basic.def.odr]}.
